% ============================================================
% 数学\数据表格\四年级.tex
% 本文件由 main.tex 通过 \input 引入，请勿单独编译
% ============================================================


\section{解答：世界上陆地面积最大的四个国家}

\vspace{0.4cm}

\noindent\textbf{4. 世界上陆地面积最大的四个国家的面积如下表：}

\vspace{0.4cm}

\begin{center}
\renewcommand{\arraystretch}{1.6}
\setlength{\tabcolsep}{15pt}
\begin{tabular}{|c|c|}
\hline
中\quad 国 & 九百六十万平方千米 \\
\hline
俄罗斯 & 一千七百零九万八千二百平方千米 \\
\hline
美\quad 国 & 九百三十七万平方千米 \\
\hline
加拿大 & 九百九十八万平方千米 \\
\hline
\end{tabular}
\end{center}

\vspace{0.6cm}

\noindent\textbf{（1）填写下表。}

\vspace{0.4cm}

\begin{center}
\renewcommand{\arraystretch}{1.6}
\setlength{\tabcolsep}{12pt}
\begin{tabular}{|c|c|c|}
\hline
国\quad 家 & 面积/平方千米 & 面积/万平方千米 \\
\hline
中\quad 国 & \textcolor{red}{9600000} & \textcolor{red}{960} \\
\hline
俄罗斯 & \textcolor{red}{17098200} & \textcolor{red}{1709.82} \\
\hline
美\quad 国 & \textcolor{red}{9370000} & \textcolor{red}{937} \\
\hline
加拿大 & \textcolor{red}{9980000} & \textcolor{red}{998} \\
\hline
\end{tabular}
\end{center}

\vspace{0.6cm}
{\small\color{gray}
\textbf{解析：}\\
首先写出各面积的阿拉伯数字形式：\\
九百六十万 $\rightarrow$ 9600000\\
一千七百零九万八千二百 $\rightarrow$ 17098200\\
九百三十七万 $\rightarrow$ 9370000\\
九百九十八万 $\rightarrow$ 9980000\\
将其改写为以“万”为单位的数，只需将小数点向左移动 4 位：\\
$9600000 = 960$ 万；
$17098200 = 1709.82$ 万；
$9370000 = 937$ 万；
$9980000 = 998$ 万。
}

\vspace{0.8cm}

%% ==============================
%% 分数的初步认识：涂色表示分数
%% ==============================
\newpage



\section{解答：体育用品商店购买问题}

\vspace{0.4cm}

\noindent\textbf{2. 一个体育用品商店，每个排球 30 元，每个足球 50 元。李老师带的钱正好能买 12 个篮球或 10 个排球。（先根据问题填写表格，再列式解答）}

\vspace{0.6cm}

% 填写表格
\begin{center}
\renewcommand{\arraystretch}{1.8}
\begin{tabular}{|c|c|c|}
\noalign{\hrule height 1pt}
\makebox[3cm][c]{品名} & \makebox[4cm][c]{单价 / (元 / 个)} & \makebox[3cm][c]{数量 / 个} \\
\hline
排球 & \textcolor{blue}{\textbf{30}} & \textcolor{blue}{\textbf{10}} \\
\hline
篮球 & \textcolor{blue}{\textbf{?}} & \textcolor{blue}{\textbf{12}} \\
\hline
足球 & \textcolor{blue}{\textbf{50}} & \textcolor{blue}{\textbf{?}} \\
\noalign{\hrule height 1pt}
\end{tabular}
\end{center}

\vspace{0.8cm}

\noindent\textbf{（1）每个篮球多少元？}

\vspace{0.2cm}

{\small\color{gray}
\textbf{解析：}\\
第一步，先求出李老师带的总钱数：排球单价 $\times$ 排球数量 $= 30 \times 10 = 300$（元）。\\
第二步，由于这笔钱正好也能买 12 个篮球，所以篮球的单价为：总钱数 $\div$ 篮球数量 $= 300 \div 12 = 25$（元）。
}

\vspace{0.2cm}
\noindent\textbf{综合算式解答：} \textbf{\textcolor{blue}{$\boldsymbol{30 \times 10 \div 12 = 25}$}} \textbf{（元）} \\
\noindent\textbf{答：}每个篮球 \textbf{\textcolor{blue}{25}} 元。

\vspace{0.6cm}

\noindent\textbf{（2）李老师带的钱能买多少个足球？}

\vspace{0.2cm}

{\small\color{gray}
\textbf{解析：}\\
同样利用前面求出的总钱数 300 元。足球的单价是 50 元/个，那么这笔钱能买足球的数量为：总钱数 $\div$ 足球单价 $= 300 \div 50 = 6$（个）。
}

\vspace{0.2cm}
\noindent\textbf{综合算式解答：} \textbf{\textcolor{blue}{$\boldsymbol{30 \times 10 \div 50 = 6}$}} \textbf{（个）} \\
\noindent\textbf{答：}李老师带的钱能买 \textbf{\textcolor{blue}{6}} 个足球。

\vspace{0.6cm}

\noindent\textbf{（3）如果买 12 个篮球和 8 个足球，一共要多少元？}

\vspace{0.2cm}

{\small\color{gray}
\textbf{解析：}\\
买 12 个篮球所需的钱就是李老师原先的总钱数（$300$ 元），或者用刚才算出的篮球单价计算：$25 \times 12 = 300$（元）。\\
买 8 个足球需要的钱为：$50 \times 8 = 400$（元）。\\
两者加起来一共要：$300 + 400 = 700$（元）。
}

\vspace{0.2cm}
\noindent\textbf{综合算式解答：} \textbf{\textcolor{blue}{$\boldsymbol{25 \times 12 + 50 \times 8 = 700}$}} \textbf{（元）} \ \textit{或直接写为：} \textbf{\textcolor{blue}{$\boldsymbol{300 + 50 \times 8 = 700}$}} \textbf{（元）}\\
\noindent\textbf{答：}一共要 \textbf{\textcolor{blue}{700}} 元。


\newpage



\section{解答：假设法解决问题——象棋与跳棋数量}

\vspace{0.4cm}

\noindent\textbf{3. 学校有象棋、跳棋共 26 副，2 人下一副象棋，6 人下一副跳棋，恰好可以供 96 人进行活动。象棋、跳棋各有多少副？（先假设两种棋的数量，再调整找出答案）}

\vspace{0.8cm}

\begin{center}
\renewcommand{\arraystretch}{1.8}
% 表格绘制：4列，包含假设过程、计算过程及结果对比
\begin{tabular}{|c|c|l|c|}
\hline
\textbf{象棋数量/副} & \textbf{跳棋数量/副} & \multicolumn{1}{c|}{\textbf{下棋总人数}} & \textbf{和 96 人比较} \\ \hline
\textcolor{blue}{13} & \textcolor{blue}{13} & \textcolor{blue}{$13 \times 2 + 13 \times 6 = 104$} & \textcolor{blue}{多 8 人} \\ \hline
\textcolor{blue}{14} & \textcolor{blue}{12} & \textcolor{blue}{$14 \times 2 + 12 \times 6 = 100$} & \textcolor{blue}{多 4 人} \\ \hline
\textcolor{blue}{15} & \textcolor{blue}{11} & \textcolor{blue}{$15 \times 2 + 11 \times 6 = 96$} & \textcolor{blue}{相等} \\ \hline
\end{tabular}
\end{center}

\vspace{0.6cm}

{\small\color{gray}
\textbf{解析：}\\
1. **第一次假设**：假设象棋和跳棋各有 $13$ 副（$26 \div 2 = 13$）。计算总人数为 $13 \times 2 + 13 \times 6 = 104$（人），比 $96$ 人多了 $8$ 人。\\
2. **分析调整方向**：因为计算出的总人数偏多，且跳棋的人均负荷（6人）高于象棋（2人），因此需要减少跳棋数量。\\
3. **逐步调整**：每减少一副跳棋并增加一副象棋，总人数减少 $6 - 2 = 4$（人）。\\
4. **得出结论**：经过两次调整，当象棋为 $15$ 副、跳棋为 $11$ 副时，总人数刚好为 $96$ 人。
}

\vspace{0.4cm}
\noindent\textbf{解答：}\\
象棋有（ \textbf{\textcolor{blue}{15}} ）副，跳棋有（ \textbf{\textcolor{blue}{11}} ）副。

\vspace{0.6cm}

{\small\color{blue!70!black}
\textbf{绘图基本原理与步骤：}\\
\textbf{1. 列表尝试法（Guess and Check）}：通过列表有序展示不同假设下的计算结果，寻找变量变化与最终数值间的规律。\\
\textbf{2. 建立逻辑关联}：表格中每一行代表一次完整的业务逻辑模拟。第一、二列为独立变量，第三列为衍生计算量，第四列为反馈修正量。\\
\textbf{3. 数据一致性}：表格内所有数据项均使用 \texttt{textcolor} 进行着色处理，既突出了动态求解过程，也保持了与题干风格的视觉统一。\\
\textbf{4. 结构化呈现}：利用 \texttt{tabular} 的水平线与垂直线构筑严谨的逻辑网格，使抽象的代数试错过程变得直观易读。
}



\newpage



\section{解答：找出等式与方程并准确连线}

\vspace{0.6cm}

\begin{center}
\begin{tikzpicture}[>=Stealth, line join=round, line cap=round, scale=1.0]
  
  % 定点坐标参数 (四行题目，采用相等的纵向间距 1.5)
  \def\yA{4.5}
  \def\yB{3.0}
  \def\yC{1.5}
  \def\yD{0.0}

  % 定义全图共用的关键图元样式
  \tikzset{
    expr/.style={font=\Large, text=black},
    box/.style={draw=black, rectangle, minimum width=2.0cm, minimum height=1.0cm, font=\Large\bfseries, text=black, line width=0.8pt},
    conn/.style={draw=red, line width=1.4pt} % 还原参考答案阅卷风格的醒目红色加粗连线
  }

  % ===================================
  % 一、 结构排版布局：三列式设计
  % ===================================

  % 1. 左侧表达式列 (利用 anchor=east 使其右侧端点严格对其到 x=0 标线上)
  \node[expr, anchor=east] (L1) at (0, \yA) {$90 - x = 30$};
  \node[expr, anchor=east] (L2) at (0, \yB) {$80 \div 4 = 20$};
  \node[expr, anchor=east] (L3) at (0, \yC) {$x - 15$};
  \node[expr, anchor=east] (L4) at (0, \yD) {$7y = 63$};

  % 2. 居中分类框列 (等式框居于 1、2 行中间，方程框居于 3、4 行中间，完美匀称)
  \node[box] (M1) at (3.5, 3.75) {等式};
  \node[box] (M2) at (3.5, 0.75) {方程};

  % 3. 右侧表达式列 (利用 anchor=west 使其左侧端点严格对其到 x=7 标线上)
  \node[expr, anchor=west] (R1) at (7, \yA) {$20 + 30 = 50$};
  \node[expr, anchor=west] (R2) at (7, \yB) {$y + 17 = 38$};
  \node[expr, anchor=west] (R3) at (7, \yC) {$36 + x < 40$};
  \node[expr, anchor=west] (R4) at (7, \yD) {$54 \div x = 9$};

  % ===================================
  % 二、 逻辑自动连线
  % (利用直接指向节点名的方式，借助 TikZ 原生算法自动在方框轮廓线生成分散连接点)
  % ===================================

  % --- 【“等式”的连线群】 ---
  % 本质要求：只要含有等号“=”的式子统统属于等式
  % 从左侧连向等式：
  \draw[conn] (L1.east) -- (M1);
  \draw[conn] (L2.east) -- (M1);
  \draw[conn] (L4.east) -- (M1);
  % 从右侧连向等式：
  \draw[conn] (R1.west) -- (M1);
  \draw[conn] (R2.west) -- (M1);
  \draw[conn] (R4.west) -- (M1);

  % --- 【“方程”的连线群】 ---
  % 本质要求：在等式的基础上，含有未知数的式子才是方程
  % 从左侧连向方程：
  \draw[conn] (L1.east) -- (M2);
  \draw[conn] (L4.east) -- (M2);
  % 从右侧连向方程：
  \draw[conn] (R2.west) -- (M2);
  \draw[conn] (R4.west) -- (M2);

\end{tikzpicture}
\end{center}

\vspace{0.6cm}

{\small\color{gray}
\textbf{解题说明：}\\
1. \textbf{概念辨析：}表示相等关系的式子叫做\textbf{等式}，所有的等式中都必须含有等号（$=$）。含有未知数的等式叫做\textbf{方程}。因此，所有的方程一定是等式，但等式不一定是方程。\\
2. \textbf{左侧判断：}\\
\quad $90-x=30$：含有未知数 $x$ 且是等式，故既连“等式”也连“方程”；\\
\quad $80\div4=20$：是等式但不含未知数，只连“等式”；\\
\quad $x-15$：不含等号，既不是等式也不是方程；\\
\quad $7y=63$：含有未知数 $y$ 且是等式，既连“等式”也连“方程”。\\
3. \textbf{右侧判断：}\\
\quad $20+30=50$：是等式但不含未知数，只连“等式”；\\
\quad $y+17=38$：含未知数且是等式，既连“等式”也连“方程”；\\
\quad $36+x<40$：含有未知数但属于不等式（含 $<$），二者皆不可连；\\
\quad $54\div x=9$：含有未知数且是等式，既连“等式”也连“方程”。
}

\vspace{0.6cm}

{\small\color{blue!70!black}
\textbf{绘图基本原理与步骤：}\\
\textbf{1. 网格化布局节点：}通过精确定义常数如 \texttt{\textbackslash def\textbackslash yA\{4.5\}} 规范四行的绝对统一高度。左侧各题目节点赋予 \texttt{anchor=east} 属性锁定同一右侧端边界，右侧题目节点相应赋予 \texttt{anchor=west} 锁定同一左侧边。以此使整体横跨三竖排的结构极端工整齐平，还原经典教辅印刷品的空间视觉秩序感。\\
\textbf{2. 智能离散连线：}在 TikZ 算法中如果指定直接向节点框体中心点相连连线（即仅使用 \texttt{-- (M1)} 而不是精准定位锚点 \texttt{-- (M1.west)}），底层原生渲染器就会自动计算出每一条带有角度跨向中心的射线与方框边界线条的物理相交点，截断并锁止在那。这一代码小幅规避了所有线段极其死板地交汇同一点的生硬状态，呈现出仿佛手工在不同下笔位自然向中心聚拢划线的动态视觉效果。\\
\textbf{3. 醒目的红勾墨迹模拟：}利用预定义常量 \texttt{conn/.style=\{draw=red, line width=1.4pt\}} 为线系批量刷上大红粗线效果，跳出题目原有静止黑白色图层，打造最高效的查漏阅卷视角。
}



\section{解答：代数式因式分解连线题}

\vspace{0.6cm}

\textbf{\LARGE 9.} \textbf{把左右两边相等的代数式用线连起来：}

\vspace{0.6cm}
\textbf{解：}

\textbf{（1）代数式因式分解恒等查验分析：}
\begin{itemize}
    \item $2a^2 - 2a$ 提取公因式 $2a$，恒等变形得到 $2a(a - 1)$；
    \item $a^2 + 6a + 9$ 是标准完全平方式，根据完全平方公式得到 $(a + 3)^2$；
    \item $4 - a^2$ 即 $2^2 - a^2$，根据平方差公式展开得到 $(2 - a)(2 + a)$；
    \item $3a^2 + 12a$ 提取公因式 $3a$，恒等变形得到 $3a(a + 4)$。
\end{itemize}

\textbf{（2）匹配连线制图：}

\begin{center}
\begin{tikzpicture}[>=Stealth, line join=round]

  % 定义左右两大正圆形的外框保护罩
  \draw[line width=0.8pt] (0,0) circle (2.2cm);
  \draw[line width=0.8pt] (8,0) circle (2.2cm);

  % 定义左侧代数式节点 (坐标以椭圆中心为基准，向左偏移排布，全量左对齐)
  \begin{scope}[every node/.style={anchor=west, inner sep=4pt}]
      \node (L1) at (-1.4, 1.2) {$2a^2 - 2a$};
      \node (L2) at (-1.4, 0.4) {$a^2 + 6a + 9$};
      \node (L3) at (-1.4, -0.4) {$4 - a^2$};
      \node (L4) at (-1.4, -1.2) {$3a^2 + 12a$};

      \node (R1) at (6.6, 1.2) {$(2 - a)(2 + a)$};
      \node (R2) at (6.6, 0.4) {$2a(a - 1)$};
      \node (R3) at (6.6, -0.4) {$(a + 3)^2$};
      \node (R4) at (6.6, -1.2) {$3a(a + 4)$};
  \end{scope}

  % 绘制穿透式等效连线 (解答区特供红色)
  \begin{scope}[draw=red!80!black, line width=1.2pt]
      \draw (L1.east) -- (R2.west);
      \draw (L2.east) -- (R3.west);
      \draw (L3.east) -- (R1.west);
      \draw (L4.east) -- (R4.west);
  \end{scope}

\end{tikzpicture}
\end{center}

\vspace{0.4cm}
{\small\color{blue!70!black}
\textbf{画图及逻辑控制思路解构：}\\
\textbf{1. 极速恒等数学推演：} 连线题的本质是多项式的因式分解死锁验证。系统在连线前于后台发起了四道极速拦截逆运算：分别命中“单项提取公因式”（$2a^2 - 2a \rightarrow 2a(a - 1)$ 和 $3a^2 + 12a \rightarrow 3a(a + 4)$）、“完全平方压铸阵”（$a^2 + 6a + 9 \rightarrow (a+3)^2$）以及极其直观的“硬式平方差暴解”（$4 - a^2 \rightarrow (2-a)(2+a)$）。确保左右每一个文本节点的神经突触实现了 $100\%$ 不可篡改的数学锁死交互，摧毁任何连错的可能。\\
\textbf{2. 坐标网格化与对称破甲布局：} 物理机层面的绘制阶段全部放弃了纯手工的游标微调；直接下令切入了横跨度达到 $8\mathrm{cm}$ 的双极制图引擎。围绕着绝对点 $(0,0)$ 和 $(8,0)$ 打下了两根半径为 $2.2\mathrm{cm}$ 的浑圆桩柱围墙，完美匹配了教辅原书的圆形圈地画风。左右的各组参数公式全部强制挂靠 \texttt{anchor=west} 左对齐吸附矩阵，以此还原卷面中那种“左侧工整，右侧锯齿错落”的文字排列真身属性。\\
\textbf{3. 穿刺连线边界锚点捕捉识别技术：} 对于这种带两边集合的交叉线网，如果直接给节点下死线，线段会直接野蛮穿断字母文本体内，这是排版大忌！为此代码启动了边框自检触雷感应，强行指令左测出线的端点绑定在 \texttt{L.east}（节点的正东方皮肤位），而降落的突击端点绑定在 \texttt{R.west}（目标节点的正西方前胸），并且挂置了 \texttt{inner sep=4pt} 的绝对弹开保护层。随后给出了 \texttt{red!80!black} 重工级粗血线执行弹射连接：可以看到四根致命的红色射线在没有任何压字重影的情况下，直接穿透了圈地圆弧墙体，犹如激光制导般精准咬合挂载在了目标算式的身侧面上，呈现出极高净度的版面洁癖美感。
}
%% ==============================
%% 第25题（补充题）：整式运算运算法则表
%% ==============================
\newpage


